\documentclass[12pt,a4paper]{article}
\usepackage{expediente}
\begin{document}
% =============================================================================
% FICHA DE DATOS DEL EXPEDIENTE --- actualizado 20 ago 2026
% =============================================================================

% --- VENDEDOR / TITULAR REGISTRAL -------------------------------------------
\newcommand{\VendedorNombres}{RICHAR DARWIN}
\newcommand{\VendedorApellidos}{CUTIRE CONDORI}
\newcommand{\VendedorNombreCompleto}{RICHAR DARWIN CUTIRE CONDORI}
\newcommand{\VendedorDNI}{42610690}
\newcommand{\VendedorEstadoCivil}{soltero, sin unión de hecho vigente}
\newcommand{\VendedorOcupacion}{\faltante{[ocupación de Richar]}}
\newcommand{\VendedorDomicilio}{\faltante{[domicilio actual de Richar]}}
\newcommand{\VendedorTelefono}{\faltante{[teléfono de Richar]}}
\newcommand{\VendedorCorreo}{\faltante{[correo]}}
\newcommand{\VendedorNacionalidad}{peruano}
\newcommand{\VendedorFechaNac}{\faltante{[fecha de nacimiento]}}
\newcommand{\VendedorDNIEmision}{\faltante{[fecha de emisión del DNI de Richar --- la pide SUNARP]}}

% --- CONVIVIENTE: NO CORRESPONDE -------------------------------------------
\newcommand{\ConvivienteNombre}{NO CORRESPONDE}
\newcommand{\ConvivienteDNI}{NO CORRESPONDE}
\newcommand{\ConvivienteDesde}{NO CORRESPONDE}
\newcommand{\ConvivienteInscrita}{no existe unión de hecho vigente}
\newcommand{\ConvivienteDomicilio}{NO CORRESPONDE}

% --- COMPRADORES (copropiedad en partes iguales) ---------------------------
\newcommand{\CompradorUnoNombre}{NICO ALVARO CUTIRE ARCE}
\newcommand{\CompradorUnoDNI}{48325017}
\newcommand{\CompradorUnoEmision}{01 de enero de 2024}
\newcommand{\CompradorUnoEstadoCivil}{soltero}
\newcommand{\CompradorUnoOcupacion}{ingeniero informático}
\newcommand{\CompradorUnoDomicilio}{Jr.\ Wiracocha, Sicuani, provincia de Canchis, Cusco, Perú}
\newcommand{\CompradorUnoTelefono}{984715108}

\newcommand{\CompradorDosNombre}{LUCY CUTIRE ARCE}
\newcommand{\CompradorDosDNI}{47478852}
\newcommand{\CompradorDosEmision}{09 de enero de 2026}
\newcommand{\CompradorDosEstadoCivil}{soltera}
\newcommand{\CompradorDosOcupacion}{contador}
\newcommand{\CompradorDosDomicilio}{Comunidad Hanocca, distrito de Layo, provincia de Canas, Cusco, Perú}
\newcommand{\CompradorDosTelefono}{953440875}

\newcommand{\CompradorNombres}{NICO ALVARO y LUCY}
\newcommand{\CompradorApellidos}{CUTIRE ARCE}
\newcommand{\CompradorNombreCompleto}{NICO ALVARO CUTIRE ARCE y LUCY CUTIRE ARCE}
\newcommand{\CompradorDNI}{48325017 y 47478852}
\newcommand{\CompradorEstadoCivil}{solteros}
\newcommand{\CompradorConyuge}{NO CORRESPONDE}
\newcommand{\CompradorConyugeDNI}{NO CORRESPONDE}
\newcommand{\CompradorOcupacion}{ingeniero informático y contador, respectivamente}
\newcommand{\CompradorDomicilio}{Sicuani (Canchis) y Layo (Canas), Cusco}
\newcommand{\CompradorTelefono}{984715108 / 953440875}
\newcommand{\CompradorCorreo}{\faltante{[correo de contacto]}}
\newcommand{\CompradorNacionalidad}{peruanos}
\newcommand{\PorcentajeCopropiedad}{50\,\% cada uno}

% --- INTERMEDIARIO / OCUPANTE DEL REMANENTE --------------------------------
\newcommand{\IntermediarioNombre}{LUCIO GIL CUTIRE PARI}
\newcommand{\IntermediarioDNI}{42467850}
\newcommand{\IntermediarioEstadoCivil}{\faltante{[estado civil de Lucio]}}
\newcommand{\IntermediarioOcupacion}{\faltante{[ocupación]}}
\newcommand{\IntermediarioDomicilio}{\faltante{[domicilio de Lucio]}}
\newcommand{\IntermediarioTelefono}{\faltante{[teléfono de Lucio]}}
\newcommand{\IntermediarioParentesco}{familiar (apellido Cutire)}

% --- PREDIO MATRIZ ----------------------------------------------------------
\newcommand{\PartidaMatriz}{\faltante{[N.º de partida electrónica SUNARP --- pendiente]}}
\newcommand{\OficinaRegistral}{Oficina Registral de Puerto Maldonado}
\newcommand{\AreaMatriz}{\faltante{[área total inscrita; las partes refieren del orden de 100 ha]}}
\newcommand{\UnidadCatastral}{\faltante{[código catastral / UC, si existe]}}
\newcommand{\NombreFundo}{\faltante{[nombre del fundo, si tiene]}}
\newcommand{\KmCarretera}{Interoceánica Sur, al este del Puente Jayave (el puente queda al oeste del predio, sobre el río; no debajo del lote)}

% --- PREDIO A INDEPENDIZAR --------------------------------------------------
\newcommand{\AreaIndependizarHa}{10{,}00}
\newcommand{\AreaIndependizarMetros}{100 000{,}00}
\newcommand{\PerimetroM}{2 200{,}58}
\newcommand{\FrenteM}{100{,}29}
\newcommand{\FondoM}{1 000{,}00}

\newcommand{\PUnoLat}{$-12{,}912452$}
\newcommand{\PUnoLon}{$-70{,}170058$}
\newcommand{\PDosLat}{$-12{,}912497$}
\newcommand{\PDosLon}{$-70{,}170981$}
\newcommand{\PTresLat}{$-12{,}903469$}
\newcommand{\PTresLon}{$-70{,}171438$}
\newcommand{\PCuatroLat}{$-12{,}903424$}
\newcommand{\PCuatroLon}{$-70{,}170515$}

\newcommand{\PUnoE}{373 065{,}00}
\newcommand{\PUnoN}{8 572 256{,}00}
\newcommand{\PDosE}{372 965{,}00}
\newcommand{\PDosN}{8 572 251{,}00}
\newcommand{\PTresE}{372 911{,}00}
\newcommand{\PTresN}{8 573 249{,}00}
\newcommand{\PCuatroE}{373 011{,}00}
\newcommand{\PCuatroN}{8 573 255{,}00}

\newcommand{\FechaPago}{15 de agosto de 2018}
\newcommand{\MontoPago}{S/\ 20 000{,}00}
\newcommand{\MontoPagoLetras}{VEINTE MIL Y 00/100 SOLES}
\newcommand{\FormaPagoHistorica}{dinero en efectivo, sin medio de pago del sistema financiero}

% Nombres genéricos --- sustituir por los reales al ir a notaría
\newcommand{\NotarioNombre}{Dr.\ JUAN CARLOS RÍOS SALINAS (nombre genérico a sustituir)}
\newcommand{\NotariaCiudad}{Puerto Maldonado}
\newcommand{\ColegiaturaAbogado}{CAL Madre de Dios N.º 0000 (genérico)}
\newcommand{\AbogadoNombre}{Abog.\ ANA MARÍA QUISPE HUAMÁN (nombre genérico a sustituir)}
\newcommand{\IngenieroNombre}{Ing.\ PEDRO ANTONIO MAMANI CONDORI (nombre genérico a sustituir)}
\newcommand{\IngenieroCIP}{CIP 000000 (genérico)}
\newcommand{\Municipalidad}{Municipalidad Distrital de Inambari}
\newcommand{\Provincia}{Tambopata}
\newcommand{\Departamento}{Madre de Dios}
\newcommand{\Distrito}{Inambari}

\newcommand{\LugarOtorgamiento}{Puerto Maldonado}
\newcommand{\DiaOtorgamiento}{\faltante{[día]}}
\newcommand{\MesOtorgamiento}{\faltante{[mes]}}
\newcommand{\AnioOtorgamiento}{2026}

\newcommand{\TestigoUno}{\faltante{[testigo 1: nombres, DNI, domicilio]}}
\newcommand{\TestigoDos}{\faltante{[testigo 2: nombres, DNI, domicilio]}}

\InicioDocumento{Proceso de compraventa del remanente a Lucio a \PrecioRemanenteLucio}{DOC-26}

\noindent Hipótesis de trabajo (no sustituye liquidación municipal ni de SUNAT): el precio de las $\sim$90 ha es \textbf{\PrecioRemanenteLucio} (\PrecioRemanenteLucioLetras), o sea \PrecioAdquisicionMatriz\ que Richar pagó en 2017 \textbf{menos} los \MontoPago\ de las 10 ha. En la minuta de Lucio eso se escribe como \textbf{precio de compraventa}, no como ``saldo'' ni ``precio suplementario'' de la escritura 962 de 2017 (Lucio no fue parte de ese acto).

\Clausula{1. Cómo queda el mapa de precios}
{\small\renewcommand{\arraystretch}{1.22}
\noindent\begin{tabularx}{\textwidth}{>{\bfseries}p{4.4cm} X}
\toprule
Acto & Precio \\
\midrule
Richar compra 100 ha (2017) & \PrecioAdquisicionMatriz\ (S/\ 27\,000 efectivo + S/\ 38\,000 Caja Arequipa) \\
Nico y Lucy --- 10 ha (2018 / formalización) & \MontoPago \\
Lucio --- remanente $\sim$90 ha & \PrecioRemanenteLucio \\
Suma de las dos ventas de Richar & S/\ 65\,000 (nominalmente recupera lo que pagó) \\
\bottomrule
\end{tabularx}}

\vspace{0.25em}
\noindent Solo vale si \textbf{Lucio realmente pagó o pagará} esos S/\ 45\,000. Si entregó otro monto, se declara el real.

\Clausula{2. Proceso de compraventa con Lucio (con este precio)}
\begin{enumerate}[leftmargin=1.15cm,itemsep=0.18em]
\item Independizar primero las 10 ha (expediente Nico/Lucy). El remanente queda en la partida de Richar.
\item Lucio: DNI, estado civil, domicilio. Si es casado o conviviente, interviene su pareja.
\item Precio en minuta: \PrecioRemanenteLucio. Forma de pago: \textbf{depósito o transferencia a la cuenta de Richar} (S/\ 45\,000 $>$ 3 UIT = S/\ 16\,500). Insertar vouchers en la escritura, igual que los de 2016 a Saturnino.
\item Predial de Richar al día (Mazuko) sobre lo que siga a su nombre.
\item Alcabala de \textbf{Lucio} en Inambari (DOC-07 adaptado al remanente).
\item Renta de segunda categoría de \textbf{Richar} ante SUNAT (o DJ de que no hay impuesto), para el notario.
\item Escritura Richar $\rightarrow$ Lucio + SID-SUNARP. Gastos: lo habitual, Lucio.
\end{enumerate}

\Clausula{3. Alcabala (la paga Lucio, no Richar)}
Base = el \textbf{mayor} entre \PrecioRemanenteLucio\ y el autoavalúo de las $\sim$90 ha, menos 10 UIT (S/\ 55\,000 en 2026). Impuesto = 3\,\% del exceso.

{\small\renewcommand{\arraystretch}{1.2}
\noindent\begin{tabularx}{\textwidth}{>{\bfseries}p{5.2cm} X}
Precio declarado & S/\ 45\,000 \ (menor que S/\ 55\,000: el precio solo no genera alcabala) \\
Si autoavalúo 90 ha $\leq$ S/\ 55\,000 & Alcabala \textbf{S/\ 0}. Igual se liquida y se saca constancia. \\
Si autoavalúo 90 ha $>$ S/\ 55\,000 & 3\,\% sobre el exceso. Ejemplo: autoavalúo S/\ 80\,000 $\rightarrow$ (80\,000 $-$ 55\,000) $\times$ 3\,\% $=$ S/\ 750. \\
\end{tabularx}}

\vspace{0.2em}
\noindent Las 90 ha tienen más chance que las 10 ha de que el municipio ponga un autoavalúo alto. El número oficial lo dice Rentas.

\Clausula{4. Impuesto a la renta / ``ganancia'' de Richar (SUNAT)}
Richar compró el \FechaAdquisicionRichar. Eso es \textbf{después} del 1.1.2004: no le alcanza la inafectación que sí tuvo Saturnino (adquirió en 1999). Tampoco es casa habitación. La venta del remanente es, en principio, \textbf{renta de segunda categoría}: 5\,\% de la ganancia.

\textbf{Ganancia} $=$ precio $-$ costo computable (el costo de 2017 se \textbf{actualiza} con el índice del MEF).

Costo de las 100 ha: \PrecioAdquisicionMatriz. Prorrateo por área (trabajo; SUNAT podría usar autoavalúo):

{\small\renewcommand{\arraystretch}{1.2}
\noindent\begin{tabularx}{\textwidth}{>{\bfseries}p{4.8cm} X}
Costo atribuible a 90 ha & $\approx$ S/\ 58\,500 \ (90\,\% de 65\,000), más actualización MEF \\
Precio Lucio & S/\ 45\,000 \\
Ganancia tentativa & \textbf{Negativa} (45\,000 $<$ 58\,500). IR $\approx$ \textbf{S/\ 0} \\
Qué lleva al notario & DJ de que \textbf{no existe impuesto por pagar} (Reglamento LIR), no el Formulario 1665 de pago. \\
\end{tabularx}}

\vspace{0.25em}
\noindent La venta de las \textbf{10 ha} a Nico/Lucy es otro cálculo: costo $\approx$ S/\ 6\,500; precio S/\ 20\,000; ganancia $\approx$ S/\ 13\,500 $\times$ 5\,\% $\approx$ \textbf{S/\ 675} (baja un poco al actualizar el costo). Eso lo paga Richar al formalizar las 10 ha, no Lucio.

\textbf{Riesgos (no ocultar):}
\begin{itemize}[leftmargin=1.1cm,itemsep=0.12em]
\item Si SUNAT solo reconoce como costo lo \textbf{bancarizado} en 2017 (S/\ 38\,000), el costo de las 90 ha baja ($\approx$ S/\ 34\,200) y S/\ 45\,000 ya dejaría ganancia. Por eso conviene archivar escritura 962 y los tres vouchers de Caja Arequipa.
\item Dos ventas en el mismo año (10 ha + 90 ha) siguen siendo segunda categoría. Una \textbf{tercera} venta en el año puede pasar a tercera (negocio).
\item Independizar y vender dos lotes no debe presentarse como loteo comercial.
\end{itemize}

\Clausula{5. Lo que no es este S/\ 45\,000}
No es donación. No es el saldo de la compraventa a Saturnino. No evita predial, notario ni SUNARP. No se ``completa'' depositando S/\ 20\,000 de 2018.

\LugarFecha
\TresFirmas{EL TITULAR}{\VendedorNombreCompleto\\DNI \VendedorDNI}{LUCIO}{\IntermediarioNombre\\DNI \IntermediarioDNI}{ENTERADOS 10 ha}{\CompradorNombreCompleto}

\end{document}
